% --- Archon physics formalization source begin ---
% archon:physics
% archon:covers PhyXMiniProblems/problem_phyx_mini_0691.lean
% archon:source-report reports/phyx_mini/problem_phyx_mini_0691.source.json

\chapter{Physics problem phyx\_mini\_0691}
\label{ch:PhyXMiniProblems_problem_phyx_mini_0691}

\paragraph{Problem source.}
\#\# Physical scenario

As some molten metal splashes, one droplet flies off to the east with initial velocity \textbackslash{}( v\_i \textbackslash{}) at angle \textbackslash{}( \textbackslash{}theta\_i \textbackslash{}) above the horizontal, and another droplet flies off to the west with the same speed at the same angle above the horizontal as shown in figure.

\#\# Question

In terms of \textbackslash{}( v\_i \textbackslash{}) and \textbackslash{}( \textbackslash{}theta\_i \textbackslash{}), find the distance between the two droplets as a function of time.

\#\# Answer choices

- A: A: \textbackslash{}( v\_i t \textbackslash{}cos \textbackslash{}theta\_i \textbackslash{})

- B: B: \textbackslash{}( 6v\_i t \textbackslash{}cos \textbackslash{}theta\_i \textbackslash{})

- C: C: \textbackslash{}( 2v\_i t \textbackslash{}cos \textbackslash{}theta\_i \textbackslash{})

- D: D: \textbackslash{}( 4v\_i t \textbackslash{}cos \textbackslash{}theta\_i \textbackslash{})

\#\# Figure caption (auxiliary; use the image as primary evidence)

The image illustrates a physics concept related to projectile motion or reflection. It depicts two identical objects approaching and leaving a surface with symmetrical paths.

Key components:

1. **Objects**: Two orange-colored objects depicted at both ends of the paths.
   
2. **Paths**: Dashed black lines show the trajectory of the objects, indicating motion towards and away from the surface.

3. **Vectors (\textbackslash{}( \textbackslash{}vec\{v\}\_i \textbackslash{}))**: Two red arrows represent velocity vectors of the objects. They are symmetrical and point along the paths of the objects.

4. **Angles (\textbackslash{}( \textbackslash{}theta\_i \textbackslash{}))**: Two angles marked between the surface and the vectors, indicating the incident and reflected angles are the same.

5. **Surface**: A horizontal line represents the surface from which the objects reflect.

This figure likely illustrates the principle of reflection, where the angle of incidence equals the angle of reflection.

\#\# Dataset metadata

Kinematics; Spatial Relation Reasoning, Predictive Reasoning

\paragraph{Recorded answer/context.}
C: C: \textbackslash{}( 2v\_i t \textbackslash{}cos \textbackslash{}theta\_i \textbackslash{})

\paragraph{Figure/image path.}
/root/proposal\_for\_physic/science-mango/phyx\_mini\_run/phyx\_data/test\_image/691.png

\paragraph{Formalization target.}
create a compiling Lean file with sorry bodies at `PhyXMiniProblems/problem\_phyx\_mini\_0691.lean`. The Lean declarations must preserve the physical quantities, dimensions or dimensional roles, figure labels, governing-law hypotheses, and final relation expressed by this problem.
Use Mathlib/Physlib names found through LeanExplore where available. If a physics API is missing, introduce faithful local abstractions rather than scalar placeholder aliases.

\begin{theorem}[Physics formalization target]
\label{thm:physics:phyx_mini_0691:target}
\lean{PhyXMiniProblems.ProblemPhyXMini0691.problem_phyx_mini_0691}
\uses{def:physics:phyx-mini-0691:phyxminiproblems-problemphyxmini0691-elapsedtimequantity, def:physics:phyx-mini-0691:phyxminiproblems-problemphyxmini0691-elapsedseconds, def:physics:phyx-mini-0691:phyxminiproblems-problemphyxmini0691-speedinmeterspersecond, def:physics:phyx-mini-0691:phyxminiproblems-problemphyxmini0691-symmetricdropletsetup, def:physics:phyx-mini-0691:phyxminiproblems-problemphyxmini0691-matchesprimaryfigure, def:physics:phyx-mini-0691:phyxminiproblems-problemphyxmini0691-hasphysicalparameters, def:physics:phyx-mini-0691:phyxminiproblems-problemphyxmini0691-satisfiesidealprojectilekinematics, def:physics:phyx-mini-0691:phyxminiproblems-problemphyxmini0691-separationinmetersat}
The assigned autoformalize agent should translate this physics problem into Lean declarations in the covered file, with theorem and lemma proof bodies written as `by sorry`.
\end{theorem}
\begin{proof}
This is an autoformalization task, not a proof task. Produce statements that can later be proved by the physics prover without weakening the physical model.
\end{proof}

% --- Archon declaration topology begin ---
\section{Declaration topology}
\begin{definition}[Elapsed Time Quantity]
  \label{def:physics:phyx-mini-0691:phyxminiproblems-problemphyxmini0691-elapsedtimequantity}
  \lean{PhyXMiniProblems.ProblemPhyXMini0691.ElapsedTimeQuantity}
  A physical elapsed-time interval
\end{definition}
\begin{proof}
  This is the declared model definition of Elapsed Time Quantity.
\end{proof}
\begin{definition}[Speed Quantity]
  \label{def:physics:phyx-mini-0691:phyxminiproblems-problemphyxmini0691-speedquantity}
  \lean{PhyXMiniProblems.ProblemPhyXMini0691.SpeedQuantity}
  A physical speed magnitude
\end{definition}
\begin{proof}
  This is the declared model definition of Speed Quantity.
\end{proof}
\begin{definition}[Acceleration Quantity]
  \label{def:physics:phyx-mini-0691:phyxminiproblems-problemphyxmini0691-accelerationquantity}
  \lean{PhyXMiniProblems.ProblemPhyXMini0691.AccelerationQuantity}
  A physical gravitational-acceleration magnitude
\end{definition}
\begin{proof}
  This is the declared model definition of Acceleration Quantity.
\end{proof}
\begin{definition}[elapsed Seconds]
  \label{def:physics:phyx-mini-0691:phyxminiproblems-problemphyxmini0691-elapsedseconds}
  \lean{PhyXMiniProblems.ProblemPhyXMini0691.elapsedSeconds}
  \uses{def:physics:phyx-mini-0691:phyxminiproblems-problemphyxmini0691-elapsedtimequantity}
  Read an elapsed-time interval in SI seconds
\end{definition}
\begin{proof}
  This is the declared model definition of elapsed Seconds.
\end{proof}
\begin{definition}[speed In Meters Per Second]
  \label{def:physics:phyx-mini-0691:phyxminiproblems-problemphyxmini0691-speedinmeterspersecond}
  \lean{PhyXMiniProblems.ProblemPhyXMini0691.speedInMetersPerSecond}
  \uses{def:physics:phyx-mini-0691:phyxminiproblems-problemphyxmini0691-speedquantity}
  Read a speed magnitude in SI metres per second
\end{definition}
\begin{proof}
  This is the declared model definition of speed In Meters Per Second.
\end{proof}
\begin{definition}[acceleration In Meters Per Second Squared]
  \label{def:physics:phyx-mini-0691:phyxminiproblems-problemphyxmini0691-accelerationinmeterspersecondsquared}
  \lean{PhyXMiniProblems.ProblemPhyXMini0691.accelerationInMetersPerSecondSquared}
  \uses{def:physics:phyx-mini-0691:phyxminiproblems-problemphyxmini0691-accelerationquantity}
  Read an acceleration magnitude in SI metres per second squared
\end{definition}
\begin{proof}
  This is the declared model definition of acceleration In Meters Per Second Squared.
\end{proof}
\begin{definition}[Droplet]
  \label{def:physics:phyx-mini-0691:phyxminiproblems-problemphyxmini0691-droplet}
  \lean{PhyXMiniProblems.ProblemPhyXMini0691.Droplet}
  The two molten-metal droplets distinguished by their launch directions
\end{definition}
\begin{proof}
  This is the declared model definition of Droplet.
\end{proof}
\begin{definition}[Horizontal Direction]
  \label{def:physics:phyx-mini-0691:phyxminiproblems-problemphyxmini0691-horizontaldirection}
  \lean{PhyXMiniProblems.ProblemPhyXMini0691.HorizontalDirection}
  The two horizontal directions named in the problem
\end{definition}
\begin{proof}
  This is the declared model definition of Horizontal Direction.
\end{proof}
\begin{definition}[Figure Quantity Label]
  \label{def:physics:phyx-mini-0691:phyxminiproblems-problemphyxmini0691-figurequantitylabel}
  \lean{PhyXMiniProblems.ProblemPhyXMini0691.FigureQuantityLabel}
  Literal physical-quantity labels repeated on the two sides of the figure
\end{definition}
\begin{proof}
  This is the declared model definition of Figure Quantity Label.
\end{proof}
\begin{definition}[sign]
  \label{def:physics:phyx-mini-0691:phyxminiproblems-problemphyxmini0691-horizontaldirection-sign}
  \lean{PhyXMiniProblems.ProblemPhyXMini0691.HorizontalDirection.sign}
  \uses{def:physics:phyx-mini-0691:phyxminiproblems-problemphyxmini0691-horizontaldirection}
  The sign of a horizontal component when east is the positive direction
\end{definition}
\begin{proof}
  This is the declared model definition of sign.
\end{proof}
\begin{definition}[Symmetric Splash Figure]
  \label{def:physics:phyx-mini-0691:phyxminiproblems-problemphyxmini0691-symmetricsplashfigure}
  \lean{PhyXMiniProblems.ProblemPhyXMini0691.SymmetricSplashFigure}
  \uses{def:physics:phyx-mini-0691:phyxminiproblems-problemphyxmini0691-droplet, def:physics:phyx-mini-0691:phyxminiproblems-problemphyxmini0691-horizontaldirection, def:physics:phyx-mini-0691:phyxminiproblems-problemphyxmini0691-figurequantitylabel}
  Qualitative evidence transcribed from the supplied bitmap. The figure shows two orange droplets, two red outward velocity arrows labelled v\_i, two angles labelled θ\_i, mirror-image dashed trajectories, and a horizontal reference surface meeting at the common launch point
\end{definition}
\begin{proof}
  This is the declared model definition of Symmetric Splash Figure.
\end{proof}
\begin{definition}[Symmetric Droplet Setup]
  \label{def:physics:phyx-mini-0691:phyxminiproblems-problemphyxmini0691-symmetricdropletsetup}
  \lean{PhyXMiniProblems.ProblemPhyXMini0691.SymmetricDropletSetup}
  \uses{def:physics:phyx-mini-0691:phyxminiproblems-problemphyxmini0691-elapsedtimequantity, def:physics:phyx-mini-0691:phyxminiproblems-problemphyxmini0691-speedquantity, def:physics:phyx-mini-0691:phyxminiproblems-problemphyxmini0691-accelerationquantity, def:physics:phyx-mini-0691:phyxminiproblems-problemphyxmini0691-droplet, def:physics:phyx-mini-0691:phyxminiproblems-problemphyxmini0691-symmetricsplashfigure}
  Independent data for the symmetric launch. Positions are metre-coordinate readouts in the fixed two-dimensional frame. In particular, neither trajectory is defined from the requested separation formula
\end{definition}
\begin{proof}
  This is the declared model definition of Symmetric Droplet Setup.
\end{proof}
\begin{definition}[Matches Primary Figure]
  \label{def:physics:phyx-mini-0691:phyxminiproblems-problemphyxmini0691-matchesprimaryfigure}
  \lean{PhyXMiniProblems.ProblemPhyXMini0691.MatchesPrimaryFigure}
  \uses{def:physics:phyx-mini-0691:phyxminiproblems-problemphyxmini0691-symmetricdropletsetup}
  The direction assignments and literal visual evidence in the primary image
\end{definition}
\begin{proof}
  This is the declared model definition of Matches Primary Figure.
\end{proof}
\begin{definition}[Has Physical Parameters]
  \label{def:physics:phyx-mini-0691:phyxminiproblems-problemphyxmini0691-hasphysicalparameters}
  \lean{PhyXMiniProblems.ProblemPhyXMini0691.HasPhysicalParameters}
  \uses{def:physics:phyx-mini-0691:phyxminiproblems-problemphyxmini0691-speedinmeterspersecond, def:physics:phyx-mini-0691:phyxminiproblems-problemphyxmini0691-accelerationinmeterspersecondsquared, def:physics:phyx-mini-0691:phyxminiproblems-problemphyxmini0691-symmetricdropletsetup}
  Physical branch conditions: the speed and gravitational acceleration are positive, and the common elevation angle lies between horizontal and vertical
\end{definition}
\begin{proof}
  This is the declared model definition of Has Physical Parameters.
\end{proof}
\begin{definition}[Satisfies Ideal Projectile Kinematics]
  \label{def:physics:phyx-mini-0691:phyxminiproblems-problemphyxmini0691-satisfiesidealprojectilekinematics}
  \lean{PhyXMiniProblems.ProblemPhyXMini0691.SatisfiesIdealProjectileKinematics}
  \uses{def:physics:phyx-mini-0691:phyxminiproblems-problemphyxmini0691-elapsedseconds, def:physics:phyx-mini-0691:phyxminiproblems-problemphyxmini0691-speedinmeterspersecond, def:physics:phyx-mini-0691:phyxminiproblems-problemphyxmini0691-accelerationinmeterspersecondsquared, def:physics:phyx-mini-0691:phyxminiproblems-problemphyxmini0691-symmetricdropletsetup}
  Ideal projectile motion in a uniform downward gravitational field, with no air resistance. Coordinate 0 points east and coordinate 1 points upward. The horizontal component has the sign supplied by the figure direction; both droplets share the same vertical component and acceleration. This governing law contains no distance-between-droplets formula
\end{definition}
\begin{proof}
  This is the declared model definition of Satisfies Ideal Projectile Kinematics.
\end{proof}
\begin{definition}[separation In Meters At]
  \label{def:physics:phyx-mini-0691:phyxminiproblems-problemphyxmini0691-separationinmetersat}
  \lean{PhyXMiniProblems.ProblemPhyXMini0691.separationInMetersAt}
  \uses{def:physics:phyx-mini-0691:phyxminiproblems-problemphyxmini0691-elapsedtimequantity, def:physics:phyx-mini-0691:phyxminiproblems-problemphyxmini0691-symmetricdropletsetup}
  Euclidean distance between the two droplet positions, read in metres
\end{definition}
\begin{proof}
  This is the declared model definition of separation In Meters At.
\end{proof}
\begin{definition}[Answer Choice]
  \label{def:physics:phyx-mini-0691:phyxminiproblems-problemphyxmini0691-answerchoice}
  \lean{PhyXMiniProblems.ProblemPhyXMini0691.AnswerChoice}
  Labels of the four multiple-choice expressions
\end{definition}
\begin{proof}
  This is the declared model definition of Answer Choice.
\end{proof}
\begin{definition}[coefficient]
  \label{def:physics:phyx-mini-0691:phyxminiproblems-problemphyxmini0691-answerchoice-coefficient}
  \lean{PhyXMiniProblems.ProblemPhyXMini0691.AnswerChoice.coefficient}
  \uses{def:physics:phyx-mini-0691:phyxminiproblems-problemphyxmini0691-answerchoice}
  Numerical coefficient multiplying v\_i * t * cos θ\_i in each choice
\end{definition}
\begin{proof}
  This is the declared model definition of coefficient.
\end{proof}
\begin{definition}[recorded Answer Choice]
  \label{def:physics:phyx-mini-0691:phyxminiproblems-problemphyxmini0691-recordedanswerchoice}
  \lean{PhyXMiniProblems.ProblemPhyXMini0691.recordedAnswerChoice}
  \uses{def:physics:phyx-mini-0691:phyxminiproblems-problemphyxmini0691-answerchoice}
  The answer label recorded in the supplied dataset metadata
\end{definition}
\begin{proof}
  This is the declared model definition of recorded Answer Choice.
\end{proof}
% --- Archon declaration topology end ---
% --- Archon physics formalization source end ---
